\begin{abstract}
Top-performing models on standard benchmarks often fail dramatically on held-out test sets --- ImageNet classifiers drop 11--15\% on ImageNet-V2, and BERT's 84\% MNLI accuracy collapses on HANS --- yet we lack predictive metrics to identify saturated benchmarks before investing in leaderboard climbing. We propose DNSI (Difficulty-Normalized Saturation Index), which measures benchmark saturation via the entropy of improvement patterns in SOTA histories, normalized by task difficulty. Our key insight is that saturated benchmarks exhibit compressed improvement entropy: remaining gains are narrow, incremental optimizations rather than diverse innovations. In a pilot study across four benchmarks with published generalization gaps, DNSI correlates with gap magnitude ($R = -0.95$, though $n=4$ yields wide confidence intervals), shows promise as a temporal predictor ($R^2 = 0.35$, with caveats noted below), and exhibits consistent direction across vision and NLP domains. DNSI enables researchers to assess benchmark health using only historical SOTA records, without constructing expensive held-out test sets --- helping redirect effort toward benchmarks that still yield generalizable progress.
\end{abstract}
